\chapter{Cardiac Arrest}

\section*{Definition and Overview}
Cardiac arrest is a medical emergency in which the heart suddenly stops pumping blood effectively, so the brain and other organs are starved of oxygen. A person in cardiac arrest collapses, becomes unresponsive, and stops breathing normally, often with only occasional gasps. Without immediate help, irreversible brain damage and death occur within minutes. The chance of survival improves dramatically when bystanders start cardiopulmonary resuscitation (CPR) immediately and a defibrillator is used early. Cardiac arrest is different from a heart attack, which is a blockage in a coronary artery; a heart attack can, however, trigger a cardiac arrest.

\section*{Epidemiology}
Out-of-hospital cardiac arrest is a major public health problem, affecting roughly 1 in 2000 adults each year in high-income countries, which translates to tens of thousands of events annually in Australia alone. It is more common in men and in older people, and most events occur in the home rather than in public places. The underlying cause is usually heart disease, particularly coronary artery disease. Survival varies widely by region and setting, but is generally much higher when a witness is present, CPR is started quickly, and a defibrillator is used early. In-hospital cardiac arrest carries a better outlook because of immediate access to monitoring and specialised resuscitation teams.

\section*{Aetiology and Pathophysiology}
Cardiac arrest usually results from a sudden disturbance of the heart's electrical rhythm. The commonest initial rhythm in adults is ventricular fibrillation, in which the lower chambers quiver chaotically instead of pumping blood. Other causes include:
\begin{itemize}
    \item \textbf{Coronary artery disease:} heart attacks and severe blockages that destabilise the heart's electrical system
    \item \textbf{Cardiomyopathy:} thickened, weakened, or rigid heart muscle that can generate dangerous rhythms
    \item \textbf{Valve disease and heart failure:} structural and functional abnormalities that place stress on the heart
    \item \textbf{Electrical (arrhythmic) syndromes:} inherited conditions such as long QT syndrome and Brugada syndrome that predispose to dangerous rhythms in otherwise healthy young people
    \item \textbf{Other causes:} drowning, severe blood loss, drug overdose, low blood oxygen, dangerously low or high blood potassium, pulmonary embolism, and severe hypothermia
\end{itemize}
In each case, the effective pumping action of the heart is lost, blood flow to the brain stops, and consciousness is lost within seconds.

\section*{Clinical Features}
Cardiac arrest is recognised by its sudden, dramatic presentation:
\begin{itemize}
    \item Sudden collapse and loss of responsiveness
    \item Not breathing, or only gasping (agonal breathing)
    \item No pulse
    \item Skin that is pale, grey, or blue
    \item Some people experience warning symptoms in the hour beforehand, such as chest pain, palpitations, dizziness, or breathlessness, and seek help in time
\end{itemize}

\section*{Differential Diagnosis}
When a person collapses and is unresponsive, several conditions must be distinguished from cardiac arrest:
\begin{itemize}
    \item Simple fainting (syncope), in which breathing and a pulse continue
    \item Seizures or fits, which often recover with breathing present
    \item Stroke causing sudden collapse and unconsciousness
    \item Severe hypoglycaemia (very low blood sugar) in people with diabetes
    \item Drug or alcohol overdose causing profound unconsciousness
    \item Choking, where a blocked airway stops breathing
\end{itemize}
In any collapsed person who is unresponsive and not breathing normally, assume cardiac arrest and start CPR without delay.

\section*{When to Seek Medical Care}
Cardiac arrest is itself a 000 emergency. Anyone who collapses and does not respond should be treated as being in cardiac arrest.
\textbf{Call 000 (or local emergency services) for:}
\begin{itemize}
    \item A person who has collapsed and is unresponsive
    \item Not breathing or only gasping
    \item Sudden severe chest pain, collapse, or loss of consciousness
    \item While waiting for help, start CPR and use an automated external defibrillator (AED) if one is available
\end{itemize}

\section*{Diagnosis and Investigations}
Cardiac arrest is diagnosed clinically, at the bedside, by unresponsiveness and abnormal or absent breathing. Diagnosis during resuscitation and the investigations that follow focus on finding and treating the cause:
\begin{itemize}
    \item \textbf{Continuous cardiac monitoring:} during resuscitation, to detect shockable rhythms such as ventricular fibrillation
    \item \textbf{Capnography:} measures exhaled carbon dioxide to confirm that airway management and chest compressions are effective
    \item \textbf{Blood tests:} electrolytes, troponin, glucose, and blood gases to identify reversible causes
    \item \textbf{Electrocardiogram (ECG):} after stabilisation, to look for heart attack patterns or inherited rhythm syndromes
    \item \textbf{Echocardiogram:} to assess heart function and structural damage
    \item \textbf{Imaging:} CT of the head or chest where a stroke, pulmonary embolism, or aortic problem is suspected
    \item \textbf{In survivors without an obvious cause:} assessment for inherited cardiac conditions, including screening of close family members
\end{itemize}

\section*{Management}
The treatment of cardiac arrest is time-critical and follows a standard sequence:
\begin{itemize}
    \item \textbf{Call for help:} dial 000 and ask for the ambulance service; put the phone on speaker so the dispatcher can guide CPR
    \item \textbf{CPR:} push hard and fast in the centre of the chest, at about 100--120 compressions per minute, alternating with rescue breaths if trained; minimise pauses
    \item \textbf{Defibrillation:} apply an AED as soon as possible and follow its voice prompts; each shock increases the chance of restoring a normal rhythm
    \item \textbf{Advanced life support:} paramedics and hospital teams use airway devices, drugs to support the heart and rhythm, and treat reversible causes
    \item \textbf{Care after return of circulation:} targeted temperature management, treatment of the underlying cause, and management in an intensive care unit
    \item \textbf{Rehabilitation:} survivors may need physical, cognitive, and psychological support, and an implantable defibrillator where arrhythmia risk remains high
\end{itemize}

\section*{Complications}
Even with successful resuscitation, cardiac arrest often leaves after-effects:
\begin{itemize}
    \item Brain injury from oxygen deprivation, ranging from subtle memory problems to permanent disability
    \item Fractures of the ribs or breastbone from chest compressions
    \item Lung damage, including aspiration of stomach contents into the lungs
    \item Kidney, liver, or other organ failure from low blood flow during the arrest
    \item Post-cardiac-arrest psychological distress, including anxiety, depression, and post-traumatic stress disorder
    \item Recurrence of cardiac arrest if the underlying cause is not identified and corrected
\end{itemize}

\section*{Prognosis and Outcomes}
The outlook depends strongly on how quickly treatment begins. Survival is far more likely when CPR is started within minutes and a shock is delivered promptly; for every minute without CPR and defibrillation, the chance of survival falls by roughly 10\%. In many out-of-hospital arrests the outcome remains poor, but bystander CPR and public-access defibrillators improve survival several-fold. Among survivors, recovery of brain function varies widely: some people return to a full and active life, while others have lasting disability. In-hospital arrests and those with a shockable rhythm generally carry a better outlook.

\section*{Prevention and Self-Care}
\begin{itemize}
    \item Learn CPR and how to use an AED, and refresh your skills regularly
    \item Recognise the early warning signs of a heart attack, such as chest discomfort, and seek help immediately
    \item Treat high blood pressure, high cholesterol, and diabetes under medical supervision
    \item Do not smoke, and limit alcohol intake
    \item Follow specialist advice for inherited heart conditions, including screening of family members
    \item If you have an implantable defibrillator, keep follow-up appointments and report any shocks or symptoms promptly
\end{itemize}

\section*{Sources and Further Reading}
The following sources provide authoritative, up-to-date information on cardiac arrest, resuscitation, and prevention.
\begin{itemize}
    \item St John Ambulance Australia. \textit{CPR and defibrillation training.} https://www.stjohn.org.au/
    \item NHS. \textit{What to do in a cardiac arrest (CPR and AED).} https://www.nhs.uk/conditions/first-aid/
    \item American Heart Association. \textit{About cardiac arrest and CPR.} https://www.heart.org/
    \item Resuscitation Council UK. \textit{Cardiac arrest and cardiopulmonary resuscitation guidance.} https://www.resus.org.uk/
    \item Mayo Clinic. \textit{Sudden cardiac arrest: symptoms and causes.} https://www.mayoclinic.org/diseases-conditions/sudden-cardiac-arrest/symptoms-causes/syc-20352734
    \item MedlinePlus. \textit{Cardiac arrest.} https://medlineplus.gov/cardiacarrest.html
\end{itemize}
